\GalSourcePageStart{062}{L0061}{33}{33}

\begin{center}
{\Large\bfseries MÉMOIRE\par}
\vspace{0.35em}
{\small\bfseries SUR LES\par}
\vspace{0.35em}
{\large\bfseries CONDITIONS DE RÉSOLUBILITÉ DES ÉQUATIONS PAR RADICAUX\textsuperscript{(*)}.\par}
\vspace{0.8em}
\rule{0.12\linewidth}{0.4pt}
\end{center}

Le Mémoire ci-joint\textsuperscript{(2)} est extrait d'un Ouvrage que j'ai eu
l'honneur de présenter à l'Académie il y a un an. Cet Ouvrage
n'ayant pas été compris, les propositions qu'il renferme ayant été
révoquées en doute, j'ai dû me contenter de donner, sous forme
synthétique, les principes généraux et une \emph{seule} application de
ma théorie. Je supplie mes juges de lire du moins avec attention
ce peu de pages.

On trouvera ici une \emph{condition} générale à laquelle \emph{satisfait
toute équation soluble par radicaux}, et qui réciproquement
assure leur résolubilité. On en fait l'application seulement aux

\GalHistoricalNote{\textsuperscript{(*)} Ce Mémoire et le suivant ont été retrouvés dans les papiers de Galois et
publiés pour la première fois en 1846 par Liouville, qui les avait fait précéder de
la note suivante:

« En insérant dans leur Recueil la lettre qu'on vient de lire, les éditeurs de la
\emph{Revue encyclopédique} annonçaient qu'ils publieraient prochainement les manuscrits
laissés par Galois. Mais \GalControlDefect{cette}{W09-CD001} promesse n'a pas été tenue. M.~Auguste Chevalier
avait cependant préparé le travail. Il nous a remis et l'on trouvera dans
les feuilles qui vont suivre:

» 1\textsuperscript{o} Un Mémoire entier sur les conditions de résolubilité des équations par
radicaux, avec l'application aux équations de degré premier;

» 2\textsuperscript{o} Un fragment d'un second Mémoire où Galois traite de la théorie générale
des équations qu'il nomme \emph{primitives}.

» Nous avons conservé la plupart des notes que M.~Auguste Chevalier avait
jointes aux Mémoires dont nous venons de parler. Ces notes sont toutes marquées
des initiales A.~Ch. Les notes non signées sont de Galois lui-même.

» Nous compléterons cette publication par quelques autres morceaux extraits
des papiers de Galois, et qui, sans avoir une grande importance, pourront cependant
encore être lus avec intérêt par les géomètres. »

Les extraits dont parle Liouville dans la dernière phrase de cette note n'ont
jamais été publiés.

\textsuperscript{(2)} J'ai jugé convenable de placer en tête de ce Mémoire la préface qu'on va
lire, bien que je l'aie trouvée biffée dans le manuscrit. \hfill \textsc{(A.~Ch.)}

\hfill E.~G.}
